\begin{corollary}[6.4.2]
\label{IV.2.s6.4.2-fr}
Sean $X$, $Y$ dos preesquemas localmente noetherianos, $f:X\to Y$ un
morfismo plano y $\mathscr E$ un $\mathcal O_Y$-módulo coherente. Supongamos
que, para todo $y\in Y$, el preesquema $f^{-1}(y)$ posee la propiedad $(S_k)$.
Entonces, para que $f^*(\mathscr E)$ satisfaga $(S_k)$ en un punto $x$, es
necesario y suficiente que $\mathscr E$ posea la propiedad $(S_k)$ en el punto
$f(x)$.
\end{corollary}

\begin{remark}[6.4.3]
\label{IV.2.s6.4.3-fr}
Sean $A$, $B$ dos anillos locales noetherianos, $k$ el cuerpo residual de
$A$, $\varphi:A\to B$ un homomorfismo local y $M$ un $B$-módulo de tipo finito
que sea un $A$-módulo plano, y supongamos además que el $A$-módulo $A$ y el
$(B\otimes_Ak)$-módulo $M\otimes_Ak$ poseen la propiedad $(S_n)$; ¿puede
concluirse entonces que el $B$-módulo $M$ posee la propiedad $(S_n)$? Lo
ignoramos incluso cuando $n=1$, $M=B$ y $B=\widehat A$; dicho de otro modo,
ignoramos si, en general, el \emph{completado} de un anillo local noetheriano
que satisface $(S_n)$ satisface también $(S_n)$, incluso para $n=1$. Poniendo
$Y=\operatorname{Spec}(A)$, $X=\operatorname{Spec}(B)$ y
$\mathscr F=\widetilde M$, bastaría, en virtud de (6.4.1, (ii)), demostrar que,
para \emph{todo} $y\in Y$, $\mathscr F_y$ posee la propiedad $(S_n)$
(y no solamente cuando $y$ es el punto cerrado de $Y$); bastaría también
demostrar que el conjunto $U$ de los $x\in X$ tales que
$\mathscr F_x\otimes_{\mathcal O_{f(x)}}\boldsymbol{k}(f(x))$ satisface $(S_n)$ es \emph{abierto}
en $X$ (puesto que contiene el punto cerrado de $X$ por hipótesis).
Demostraremos más adelante (12.1.4) que $U$ es abierto cuando $B$ es un anillo
local de una $A$-álgebra de tipo finito y $M=B$. Por otra parte, para
$B=\widehat A$ y $M=B$, hemos visto en (6.3.8) que los preesquemas
$f^{-1}(y)$ son preesquemas de Cohen--Macaulay (es decir, satisfacen
$(S_n)$ para \emph{todo} $n$) cuando se supone que $A$ es cociente de un
anillo local de Cohen--Macaulay. Se concluye que, cuando $A$ es cociente de
un anillo local de Cohen--Macaulay (o, más generalmente, de un anillo local
que satisface las hipótesis de (6.3.8)), \emph{para que $A$ satisfaga $(S_n)$,
es necesario y suficiente que su completado $\widehat A$ satisfaga $(S_n)$}.
Quedaría por ver si esta propiedad subsiste sin restricción sobre $A$.
\end{remark}
